% !PRINTED 554
\weberpage{554}

%<<CONT
\weq{3}{\begin{aligned}
\frac{Z^0(t)}{t} &= \frac{\pi}{2\,\beta} + \frac{\gamma^0}{\sqrt{t}}, \\
\frac{Z^1(t)}{t} &= \frac{\pi}{2\,\beta} + \frac{\gamma^1}{\sqrt{t}},
\end{aligned}}

worin $\gamma^0$ und $\gamma^1$ für $t=\infty$ endlich bleiben. Der Faktor $\frac{1}{2}$ bei
$\pi/2\,\beta$ kommt daher, daß hier die Gittermaschen Rechtecke vom Inhalt 2 sind. Dies ist
aber in Übereinstimmung mit der Formel § 146, (2).

Aus (3) ergibt sich

\weq{4}{Z^0(t) - Z^1(t) = \gamma\sqrt{t},}

worin $\gamma$ gleichfalls endlich bleibt.

Nun können wir die in der Formel (12), § 145 vorkommende Summe
\[
\sigma = \sum \frac{(-1)^x}{[(x-\alpha y)^2 + \beta^2 y^2]^s}
\]
so schreiben:
\[
\sigma = \frac{\varepsilon_1}{\mu_1^s} + \frac{\varepsilon_2}{\mu_2^s} + \frac{\varepsilon_3}{\mu_3^s} + \cdots,
\]
worin die $\varepsilon_1$, $\varepsilon_2$, $\varepsilon_3$, \dots\ nur die Werte $\pm 1$ haben und

\weq{5}{\begin{aligned}
\varepsilon_1 + \varepsilon_2 + \varepsilon_3 + \cdots + \varepsilon_n &= Z^0(n) - Z^1(n), \\
&= \gamma \sqrt{n}.
\end{aligned}}

Damit sind die Voraussetzungen unseres Satzes § 146 erfüllt.

Es ist also $\sigma$ für $s > \frac{1}{2}$ eine stetige Funktion von $s$, also
insbesondere auch für $s=1$, und wir erhalten aus § 145, (12):
\weq{6}{\begin{aligned}
-\frac{2\pi}{\beta}\log f(\omega_1)\,f(\omega_2) &= \sum \frac{(-1)^{x+y}}{(x-\alpha y)^2 + \beta^2 y^2}, \\
-\frac{2\pi}{\beta}\log f_1(\omega_1)\,f_1(\omega_2) &= \sum \frac{(-1)^{x}}{(x-\alpha y)^2 + \beta^2 y^2}, \\
-\frac{2\pi}{\beta}\log f_2(\omega_1)\,f_2(\omega_2) &= \sum \frac{(-1)^{y}}{(x-\alpha y)^2 + \beta^2 y^2}.
\end{aligned}}

Man kann diese Formeln auch noch anders darstellen. Wir führen die Bezeichnung ein:

\weq{7}{\begin{aligned}
\sum \frac{1}{(x-\alpha y)^2 + \beta^2 y^2} &= S_{00} : x \text{ gerade}, \ y \text{ gerade}, \\
&= S_{01} : x \text{ gerade}, \ y \text{ ungerade}, \\
&= S_{10} : x \text{ ungerade}, \ y \text{ gerade}, \\
&= S_{11} : x \text{ ungerade}, \ y \text{ ungerade}.
\end{aligned}}
% [sic] In(5) wird die Anzahl n der Folgenglieder mit der Schranke n der quadratischen Form vermischt; gedruckte Argumente unveraendert.
